% =========================================================
%  frontmatter/petunjuk.tex
%  Petunjuk pemakaian modul.
% =========================================================

\section*{Petunjuk Pemakaian}
\addcontentsline{toc}{section}{Petunjuk Pemakaian}

\subsection*{Untuk siapa modul ini}

Modul ini dipakai mahasiswa Program Studi \NamaProdi\ yang menempuh mata
kuliah \NamaMK\ dan baru pertama kali membangun permainan. Tidak ada
prasyarat pemrograman permainan; yang dibutuhkan hanya dasar Algoritma dan
Pemrograman: variabel, percabangan, perulangan, fungsi, kelas, dan cara
membaca pesan kesalahan.

Anda tidak perlu menghafal isi modul. Yang perlu Anda lakukan adalah
mengetik, menjalankan, mengamati layar, lalu memperbaiki bila hasilnya
belum sesuai. Setiap sesi berakhir pada program yang berjalan.

\subsection*{Urutan sesi dan perkiraan waktu}

Setiap sesi dirancang untuk satu pertemuan praktikum, sekitar 100--150 menit
termasuk mencoba dan mengisi lembar kerja. Kerjakan berurutan; sesi
berikutnya selalu melanjutkan berkas sesi sebelumnya.

{\small
\renewcommand{\arraystretch}{1.25}
\begin{tabularx}{\linewidth}{@{}C{1.1cm} L{6.4cm} Y C{1.9cm}@{}}
\toprule
\textbf{Sesi} & \textbf{Judul} & \textbf{Yang ditambahkan} & \textbf{Waktu} \\
\midrule
01 & Lingkungan Kerja dan Jendela Pertama &
     Perkakas, struktur folder, jendela berjudul & 100 menit \\
02 & Siklus Permainan dan Pengendali Waktu &
     \istilah{Game loop}, \istilah{event}, \istilah{Clock}, FPS & 110 menit \\
03 & Pemain Bergerak dengan Vektor dan Delta Time &
     \code{Vector2}, masukan papan tombol, \istilah{delta time} & 120 menit \\
04 & Gravitasi, Lompatan, dan Pantulan &
     Integrasi Euler, lompat, pantulan dinding & 130 menit \\
05 & Sprite, Artefak, Peti, dan Pintu &
     \istilah{Sprite}, \istilah{Group}, \istilah{collision}, skor, nyawa & 140 menit \\
06 & Level Lorong dan Penjaga yang Mengejar &
     Level tampak atas, penjaga reaktif, garis arah & 130 menit \\
07 & Mesin Keadaan Penjaga &
     \istilah{Finite State Machine}: berdiam, berpatroli, mengejar & 140 menit \\
08 & Grid dan Pencarian Jalur A* &
     Grid, \code{g}, \code{h}, \code{f = g + h}, animasi jalur & 150 menit \\
09 & Integrasi FSM dan A*, Dua Level Berurutan &
     \istilah{repath}, alur dua level, tiga skenario uji & 150 menit \\
\bottomrule
\end{tabularx}
}

\subsection*{Cara menjalankan checkpoint Python}

Setiap sesi meninggalkan satu folder proyek utuh pada \file{code/sesi-NN}.
Folder itu berisi \file{main.py} dan paket \file{gudang_artefak}. Untuk
menjalankannya, buka terminal pada folder sesi tersebut, lalu ketik:

\begin{lstlisting}[language={}, numbers=none, caption={Menjalankan checkpoint}]
cd code/sesi-05
python main.py
\end{lstlisting}

Bila muncul pesan bahwa Pygame tidak ditemukan, pasang lebih dahulu dengan
\code{python -m pip install pygame}. Modul ini memakai Python 3.8 atau lebih
baru.

Untuk menguji tanpa duduk lama di depan layar, tambahkan pilihan
\code{--frames 120}. Program akan berhenti sendiri setelah 120 bingkai.
Pilihan \code{--level 2} langsung membuka Level Lorong, dan
\code{--scenario 1}, \code{2}, atau \code{3} memilih peta Lorong.

\subsection*{Bila sesi sebelumnya gagal}

Jangan mengulang seluruh modul. Buka folder checkpoint sesi terakhir yang
berhasil Anda jalankan, salin ke folder kerja Anda, lalu kerjakan sesi
berikutnya dari salinan itu. Cara ini sah dan memang disediakan: yang
dinilai adalah kemampuan Anda menjelaskan dan mengubah program, bukan
mengetik ulang berkas yang sudah benar.

Bila program berhenti dengan pesan kesalahan, bacalah dari baris terakhir ke
atas. Baris terakhir menyebut jenis kesalahan; baris di atasnya menyebut
berkas dan nomor baris penyebabnya.

\subsection*{Hubungan dengan buku ajar dan RPS}

Buku ajar menjelaskan mengapa sesuatu bekerja: definisi, rumus, dan
peringatan. Modul ini menjelaskan bagaimana membangunnya: berkas apa yang
dibuat, fungsi apa yang ditulis, dan nilai apa yang diisi. Setiap sesi
menyebut nomor bab yang menjadi sumber konsepnya, sehingga Anda dapat
berpindah antara modul dan buku bila memerlukan penjelasan tambahan.

RPS mengatur capaian dan asesmen. Sembilan sesi modul ini menuntaskan
\ModulSubCPMK. Tugas praktikum pada RPS dinilai dari checkpoint yang Anda
hasilkan pada sesi yang bersangkutan.
